\chapter{Indian Myths}
This chapter focuses on medical beliefs particularly associated with parts of Indian society. These beliefs are not held by all Indians or Hindus; many are regional, caste- or community-specific, family traditions, or practices that predate or extend beyond Hinduism. Cultural significance and medical evidence are different questions, and respectful discussion should distinguish the two.

\newcommand{\indianentry}[5]{%
\item[] \textbf{#1.} \textbf{Myth:} #2\newline
\textbf{Status:} #3.\newline
\textbf{Fact:} #4 \cite{#5}
}
\begin{enumerate}[leftmargin=0pt,labelwidth=0pt,labelsep=0pt,itemindent=0pt]
\indianentry{IND-001}{Menstruating women are physically ``impure'' and can contaminate food.}{False}{Menstruation is a normal biological process, not medical uncleanliness. Menstrual status does not make food poisonous or contaminated; food safety depends on hygiene, water, storage, temperature, and microorganisms. Stigma can instead reduce access to washing facilities, pain relief, education, and care.}{ref67}
\indianentry{IND-002}{A menstruating woman touching pickles makes them spoil.}{False}{Menstruation has no effect on pickle preservation. Spoilage depends on microorganisms, moisture, salt, acidity, oxygen, temperature, and container hygiene. Cultural rules may be observed voluntarily, but they are not microbiology.}{ref67}
\indianentry{IND-003}{Menstruating women should not enter the kitchen because their presence affects food.}{False}{There is no biological mechanism by which menstruation changes food or cooking. The relevant risks are ordinary burns and food contamination, not menstrual status. Kitchen access should be based on safety and personal choice.}{ref67}
\indianentry{IND-004}{Menstruating girls should not wash their hair or bathe.}{False}{Bathing and washing hair during menstruation are safe and may support comfort and hygiene. Menstruation does not make ordinary bathing dangerous. Restrictions that prevent washing can increase discomfort and stigma.}{ref67}
\indianentry{IND-005}{Menstrual blood is ``dirty blood'' or toxins being expelled from the body.}{False}{Menstrual fluid contains blood, cervical mucus, vaginal secretions, and shed endometrial tissue. It is not a monthly detoxification process. Very heavy, prolonged, or unusually painful bleeding deserves medical assessment.}{ref67}
\indianentry{IND-006}{Girls should avoid curd, tamarind, lemon, or other sour foods during menstruation.}{False}{Sour foods do not inherently stop or disturb menstruation. Individual intolerance or reflux may justify avoiding a particular food, but there is no universal menstrual restriction. A balanced diet and adequate fluids remain appropriate.}{ref67}
\indianentry{IND-007}{Eating papaya during pregnancy necessarily causes miscarriage.}{False}{Normal dietary consumption of ripe papaya has not been shown to cause miscarriage. Unripe or semi-ripe papaya is a different exposure and is sometimes avoided as a precaution. Bleeding or pain in pregnancy requires clinical assessment, not a fruit-based remedy.}{ref13}
\indianentry{IND-008}{Eating saffron during pregnancy makes the baby fair.}{False}{Fetal skin pigmentation is determined mainly by inherited biology and melanin pathways, not by the color of food or spices. Saffron cannot make a baby fair. Large supplemental doses should not be self-prescribed during pregnancy.}{ref13}
\indianentry{IND-009}{Drinking coconut water during pregnancy makes the baby fair.}{False}{Coconut water does not determine fetal skin color. It can be a beverage but does not replace prenatal care or balanced nutrition. People with kidney disease or fluid or potassium restrictions need individualized advice.}{ref13}
\indianentry{IND-010}{Eating dark-colored foods during pregnancy can make the baby dark.}{False}{The color of maternal food does not transfer to fetal skin. Pigmentation depends on genetics and developmental biology. This claim can reinforce colorism without providing any health benefit.}{ref13}
\indianentry{IND-011}{Eating more ghee late in pregnancy lubricates the birth canal and produces an easier delivery.}{False}{Eating ghee does not lubricate the uterus, cervix, or birth canal. Labor depends on cervical change, contractions, fetal position, pelvic anatomy, and other factors. Extra ghee adds energy but cannot prevent obstructed labor or replace skilled maternity care.}{ref13}
\indianentry{IND-012}{Pregnant women must remain indoors during an eclipse to prevent birth defects.}{False}{Solar or lunar eclipses do not cause congenital abnormalities. Birth defects have genetic, developmental, infectious, nutritional, medication-related, and other causes. Pregnant people should follow ordinary safety advice, including safe solar viewing during a solar eclipse.}{ref13}
\indianentry{IND-013}{Using a knife or scissors during an eclipse can cause a cleft lip or deformity in the baby.}{False}{There is no biological connection between cutting an object during an eclipse and fetal facial development. Cleft lip and other congenital conditions arise through complex developmental pathways. Ordinary household activities do not imprint injuries on a fetus.}{ref13}
\indianentry{IND-014}{Cutting vegetables during an eclipse can give the baby corresponding cuts or deformities.}{False}{A cut vegetable cannot transmit a cut to an unborn baby. Fetal development is not controlled by visual resemblance, thought, or activity during an eclipse. Prenatal care and nutrition are medically relevant; ritual explanations cannot diagnose congenital conditions.}{ref13}
\indianentry{IND-015}{A high, low, round, or pointed abdomen predicts the baby's sex.}{False}{Abdominal shape varies with body build, muscle tone, pregnancy number, fetal position, and gestational age. It cannot reliably identify fetal sex. Visual guesses should not replace appropriate clinical testing or prenatal care.}{ref13}
\indianentry{IND-016}{Sweet versus salty pregnancy cravings indicate whether the fetus is a boy or girl.}{False}{Cravings do not reliably identify fetal sex. They may reflect appetite, habit, nausea patterns, or changing food preferences. Persistent cravings for non-food substances should be discussed with a clinician because they may signal pica or anemia.}{ref13}
\indianentry{IND-017}{A fetal heart rate above or below a particular number tells whether the baby is a boy or girl.}{False}{Fetal heart rate changes with gestational age, activity, sleep, medicines, and temporary physiology. A single number cannot determine fetal sex. Monitoring should be interpreted by maternity professionals as an assessment of wellbeing, not as a sex test.}{ref13}
\indianentry{IND-018}{Colostrum is old, dirty, or harmful milk and should be thrown away.}{False}{Colostrum is early breast milk produced after birth and contains nutrients and immune components suited to a newborn. Its yellow color does not mean it is spoiled. Discarding it can deprive a newborn of early feeding and protection.}{ref68}
\indianentry{IND-019}{A newborn should receive honey before breast milk.}{False}{Honey should not be given to infants younger than 12 months because it can contain spores that cause infant botulism. Breast milk or an appropriate medical substitute is safer. Honey provides no special cleansing or immunity benefit.}{ref69}
\indianentry{IND-020}{A newborn should receive ghutti as a traditional first feed.}{False}{Routine prelacteal feeds are unnecessary and can interfere with early breastfeeding. Homemade or commercial mixtures may be contaminated, incorrectly concentrated, allergenic, or pharmacologically active. Newborn feeding should follow maternity and pediatric guidance.}{ref68}
\indianentry{IND-021}{Janam-ghutti cleans a newborn's stomach or intestines.}{False}{A healthy newborn gastrointestinal tract does not require ritual cleansing. The intestine is adapted to receive milk, and colostrum supports normal feeding. Unregulated mixtures can introduce microbes, sugar, alcohol, medicines, or toxic substances.}{ref68}
\indianentry{IND-022}{Putting kajal or surma in a baby's eyes improves eyesight.}{False}{Kajal and surma do not improve visual development. Some preparations contain lead or other contaminants, and application can injure or infect the eye. Redness, discharge, swelling, or poor visual response needs pediatric or ophthalmic review.}{ref79}
\indianentry{IND-023}{Kajal protects babies from nazar and therefore prevents illness.}{False}{Kajal may have cultural or cosmetic meaning, but it has no demonstrated disease-preventive effect. It cannot prevent infection or developmental conditions. Contaminated products can create lead-exposure or eye-infection risks.}{ref79}
\indianentry{IND-024}{Applying oil into a newborn's nose and ears is necessary for health.}{False}{Routine oil instillation is not needed to keep a newborn's nose or ears healthy. Oil can irritate tissue, and inserting objects can injure the ear canal or airway. Discharge, fever, or breathing difficulty should not be treated with household oil.}{ref70}
\indianentry{IND-025}{Forceful traditional infant massage makes bones stronger or changes body shape.}{False}{Gentle touch may be soothing, but forceful manipulation does not reshape bones or reliably increase bone strength. Infants have vulnerable joints, nerves, skin, and developing bones. Excessive pressure can cause bruising, fractures, dislocation, or neurologic injury.}{ref70}
\indianentry{IND-026}{Pressing a newborn girl's swollen breasts removes ``witch's milk.''}{False}{Temporary breast enlargement can occur because of maternal hormones and usually resolves without squeezing. Pressing can damage tissue and introduce infection. Redness, warmth, fever, or pus needs prompt pediatric assessment.}{ref70}
\indianentry{IND-027}{Turmeric, ash, ghee, or oil applied to the umbilical stump helps it heal.}{False}{Clean, dry cord care is generally safer than applying household substances. Powders, ash, oil, ghee, or pastes can carry microorganisms and delay recognition of infection. Spreading redness, pus, foul odor, fever, or poor feeding requires urgent care.}{ref70}
\indianentry{IND-028}{A baby's head must be repeatedly shaved to make the hair grow thicker.}{False}{Shaving cuts the visible hair shaft but does not increase the number of follicles. Hair may look darker or coarser while growing back because the ends are blunt. Shaving also carries avoidable risks of skin cuts and infection.}{ref70}
\indianentry{IND-029}{Mundan permanently improves hair density.}{False}{Ritual shaving cannot create new follicles or permanently increase hair density. Apparent thickness after shaving is a visual effect caused by uniformly short shafts. Persistent hair loss or scalp inflammation should be evaluated medically.}{ref70}
\indianentry{IND-030}{A black thread, black dot, or amulet medically protects a child from disease caused by nazar.}{False}{These objects may have cultural meaning but have no demonstrated ability to prevent infection or other disease. They should never replace vaccination, safe water, nutrition, or medical care. Tight threads can injure skin and small objects can become choking hazards.}{ref70}
\indianentry{IND-031}{Illness following nazar lagna can be diagnosed from chillies, lemon, or salt.}{False}{The appearance or behavior of such objects cannot identify the cause of an illness. Symptoms require history, examination, and appropriate testing. Rituals may provide comfort but cannot distinguish infection, poisoning, dehydration, or an emergency.}{ref30}
\indianentry{IND-032}{Burning chillies can determine whether someone has the evil eye.}{False}{Smoke and smell are affected by moisture, oils, airflow, and the material itself. They cannot diagnose disease or supernatural causation. Burning material can irritate the eyes and lungs, especially in children and people with asthma.}{ref30}
\indianentry{IND-033}{Epileptic seizures are caused by possession, ghosts, or black magic.}{False}{Epileptic seizures result from abnormal electrical activity in the brain, although other conditions can look similar. A first seizure requires medical evaluation. During a seizure, protect the person from injury, do not restrain them, and never put objects in the mouth.}{ref71}
\indianentry{IND-034}{Making someone having a seizure smell shoes can revive them.}{False}{Smelling shoes or pungent substances does not stop seizure activity or restore consciousness. It may expose the person to irritants and distract bystanders from preventing injury. Time the seizure and seek emergency help when it is prolonged, repeated, or a first event.}{ref71}
\indianentry{IND-035}{Giving an iron object or key to a person during a seizure stops it.}{False}{Holding metal does not alter the brain's electrical activity. It can cause cuts, choking, or injury if forced into the hand or mouth. Clear hazards, cushion the head, and place the person on their side after convulsions allow it.}{ref71}
\indianentry{IND-036}{Tantric or exorcism rituals can cure epilepsy.}{False}{Rituals do not treat the abnormal electrical activity that causes epilepsy and should not replace neurologic assessment or prescribed medicine. Spiritual support may coexist with care if it is safe and voluntary. Restraint, injury, starvation, or treatment cessation is dangerous.}{ref71}
\indianentry{IND-037}{Mental illness results from possession, upari hawa, black magic, or ghosts.}{False}{Mental disorders arise through interacting biological, psychological, developmental, and social factors. Cultural beliefs may influence how distress is described, but possession is not a medical explanation for all psychiatric symptoms. Stigma and coercive rituals can worsen risk and delay treatment.}{ref30}
\indianentry{IND-038}{A jhaad-phoonk ritual can medically cure psychiatric disease.}{False}{Ritual practices are not substitutes for evidence-based psychiatric treatment. Depending on the condition, effective care may include psychotherapy, medication, social support, or emergency protection. Suicidal thoughts, severe confusion, or dangerous behavior require urgent help.}{ref30}
\indianentry{IND-039}{Cow urine cures cancer.}{False}{No reliable clinical evidence establishes cow urine as a cancer cure. Cancer requires diagnosis and treatment tailored to its type and stage. Delaying oncology care for an unproven remedy can be dangerous. Natural or traditional origin does not establish effectiveness, purity, or safety.}{ref48}
\indianentry{IND-040}{Cow urine cures diabetes.}{False}{Cow urine is not an established treatment for diabetes and cannot replace glucose monitoring, nutrition care, or prescribed medicines. Uncontrolled diabetes can damage the eyes, kidneys, nerves, heart, and blood vessels. Stopping treatment can cause severe hyperglycemia or ketoacidosis.}{ref48}
\indianentry{IND-041}{Cow urine can replace medicines because it is a natural disinfectant or medicine.}{False}{Natural origin does not establish dose, purity, efficacy, or safety. A product used outside a validated indication may contain contaminants, interact with medicines, or delay effective care. Disinfection of surfaces is not the same as treatment of an infection inside the body.}{ref48}
\indianentry{IND-042}{Cow dung applied to wounds promotes healing and prevents infection.}{False}{Cow dung can contaminate an open wound with bacteria, parasites, and spores, including organisms capable of causing tetanus. It does not provide sterile wound care. Rinse wounds with clean running water, control bleeding, and obtain medical advice for deep or contaminated wounds.}{ref48}
\indianentry{IND-043}{Cow dung applied to a newborn's umbilical stump helps healing.}{False}{Applying dung to a newborn's cord can introduce infection, including tetanus. Clean, dry cord care and skilled newborn care are safer. Redness, swelling, discharge, foul odor, fever, poor feeding, or lethargy requires urgent evaluation.}{ref48}
\indianentry{IND-044}{Ganga water cannot contain disease-causing organisms because it is sacred or self-purifying.}{False}{Religious significance does not guarantee microbiological safety. Untreated surface water can contain bacteria, viruses, parasites, and chemical contaminants. Clear or flowing water is not proof that it is safe, especially for infants and people with weakened immunity.}{ref48}
\indianentry{IND-045}{Panchagavya can cure major diseases.}{False}{Broad claims that panchagavya cures major diseases are not supported by robust clinical evidence. Preparations can vary in composition and may be contaminated. Serious illness requires diagnosis and therapies shown to improve outcomes.}{ref80}
\indianentry{IND-046}{Tulsi can prevent or cure virtually every infection.}{False}{Tulsi contains biologically active compounds, but laboratory activity is not the same as proven treatment in patients. Different infections require different diagnoses, medicines, doses, and durations. Tulsi cannot replace antibiotics, antivirals, vaccination, or emergency care when indicated.}{ref80}
\indianentry{IND-047}{Haldi doodh can cure bacterial or viral infections.}{False}{Turmeric milk may be consumed as food and can be comforting, but it is not a universal antimicrobial treatment. It cannot replace indicated antibiotics, antivirals, fluids, or monitoring. Concentrated products may interact with medicines or be unsafe in pregnancy.}{ref80}
\indianentry{IND-048}{Neem can cure virtually any skin disease.}{False}{Skin disorders have different causes, including infection, allergy, inflammation, autoimmune disease, and cancer. Neem is not a universal treatment and household preparations are not standardized or sterile. A spreading rash, painful lesion, pus, fever, or non-healing sore should be assessed.}{ref80}
\indianentry{IND-049}{Karela juice can cure diabetes and eliminate the need for medication.}{False}{Evidence does not establish karela as a replacement for proven diabetes treatment. Juice preparations vary in dose and can contribute to low blood glucose with medicines. Stopping prescribed therapy can produce dangerous hyperglycemia.}{ref81}
\indianentry{IND-050}{Jamun seeds can cure diabetes permanently.}{False}{There is insufficient clinical evidence that jamun seeds permanently cure diabetes. Diabetes is managed through individualized nutrition, activity, monitoring, and medicines when indicated. A lower glucose reading does not prove that the disease has been eliminated.}{ref81}
\indianentry{IND-051}{Jaggery is sugar-free and safe in unlimited amounts for diabetes.}{False}{Jaggery contains substantial sugars and can raise blood glucose. Traditional production does not make it carbohydrate-free or harmless in unlimited amounts. The relevant issue is the total carbohydrate pattern and portion.}{ref49}
\indianentry{IND-052}{Desi ghee can be eaten without limit because it is inherently heart-protective.}{False}{Ghee is energy-dense and contains substantial saturated fat. Health effects depend on quantity, replacement foods, total dietary pattern, and individual risk factors. Traditional preparation does not make unlimited intake heart-protective.}{ref49}
\indianentry{IND-053}{Rock salt or sendha namak is harmless for people with hypertension.}{False}{Rock salt still contains substantial sodium. Substituting one salt for another without reducing total sodium may not lower blood-pressure risk. People with hypertension, heart failure, or kidney disease should follow individualized dietary advice.}{ref49}
\indianentry{IND-054}{Black salt is essentially sodium-free.}{False}{Black salt contains sodium and is not a free-use substitute for people restricting sodium. Its flavor may allow some people to use less, but that depends on the amount consumed. Natural or mineral-rich does not mean sodium-free.}{ref49}
\indianentry{IND-055}{Water stored in a copper vessel cures numerous diseases.}{False}{Copper is an essential trace element, but copper-vessel water is not a general medicine. Evidence for broad treatment claims is inadequate, and excess copper exposure can cause harm. Water safety also depends on source contamination and storage hygiene.}{ref80}
\indianentry{IND-056}{Water kept overnight in copper cures digestive and metabolic disease.}{False}{Overnight storage does not transform water into a cure. Copper exposure depends on the vessel, water chemistry, contact time, and cleaning, while excessive exposure can be harmful. Ongoing symptoms need evaluation for their actual cause.}{ref80}
\indianentry{IND-057}{Repeatedly drinking hot water “melts” body fat.}{False}{Hot water does not physically melt stored adipose tissue. Body fat changes through sustained energy balance, influenced by biology and environment. Very hot drinks can burn the mouth and esophagus.}{ref5}
\indianentry{IND-058}{Ajwain water burns abdominal fat.}{False}{Ajwain water does not selectively metabolize abdominal adipose tissue. No drink can target fat loss from one body region. Weight management should use sustainable nutrition, activity, sleep, and medical support when needed.}{ref5}
\indianentry{IND-059}{Lemon-honey water detoxifies the liver.}{False}{The liver performs its own biochemical processing; lemon and honey do not wash toxins out of it. Honey is not suitable for infants under 12 months and adds sugar. Jaundice or other liver symptoms require medical assessment.}{ref5}
\indianentry{IND-060}{Eating curd at night produces harmful mucus in everyone.}{False}{There is no general physiological rule that curd causes pathological mucus at night. Some people may have reflux, lactose-related symptoms, or intolerance, which is different from a universal mucus effect. Persistent cough or breathing symptoms need evaluation.}{ref5}
\indianentry{IND-061}{Milk and fish together produce vitiligo.}{False}{This food combination does not cause vitiligo. Vitiligo involves loss or dysfunction of pigment-producing melanocytes and is commonly associated with autoimmune mechanisms. Unnecessary avoidance can impair nutrition and reinforce stigma.}{ref82}
\indianentry{IND-062}{Milk followed by sour food causes vitiligo.}{False}{Eating sour food after milk does not cause vitiligo. Vitiligo is not an incompatibility reaction between foods and is not caused by food colors mixing in the body. New or changing depigmented patches should be assessed.}{ref82}
\indianentry{IND-063}{Vitiligo is contagious.}{False}{Vitiligo is not an infection and does not spread by touch, food, or shared clothing. It can have autoimmune and genetic contributors. Stigma may cause isolation or delayed care, while dermatologic assessment can confirm the diagnosis.}{ref82}
\indianentry{IND-064}{Leprosy results from sins committed in a previous life.}{False}{Leprosy is an infectious disease caused by Mycobacterium leprae, not a punishment or moral failure. Early diagnosis and multidrug therapy can cure the infection and reduce disability. Stigma and delayed treatment harm patients and families.}{ref74}
\indianentry{IND-065}{Leprosy spreads simply by touching someone once.}{False}{Casual contact or a single touch is not the usual route of transmission. Prolonged, close exposure to an untreated person is more relevant, and effective treatment reduces infectiousness. Suspicious skin patches, numbness, or nerve symptoms should prompt evaluation without blame.}{ref74}
\indianentry{IND-066}{Snakebite can be cured by a mantra.}{False}{Chanting does not neutralize venom. Suspected snakebite requires rapid transport, limb immobilization, and assessment for antivenom and supportive care. Do not cut, suck, burn, or apply chemicals to the bite.}{ref72}
\indianentry{IND-067}{A snake-stone or nagmani can pull venom from a snakebite.}{False}{A stone cannot extract venom that has entered tissue and the bloodstream. Attaching or cutting around it can injure the skin and waste time. Keep the person still and arrange urgent transport for professional assessment.}{ref72}
\indianentry{IND-068}{An ojha, bhopa, or tantrik can remove snake poison before hospital treatment.}{False}{Ritual intervention cannot remove systemic venom. Delay increases the risk of paralysis, bleeding, tissue injury, kidney failure, and death. The victim should be transported promptly without unnecessary walking or a tight tourniquet.}{ref72}
\indianentry{IND-069}{Turmeric, lime, or herbal paste neutralizes snake venom.}{False}{These substances do not neutralize venom circulating through the body. They can irritate the wound, introduce infection, and delay antivenom evaluation. Keep the person calm and still and obtain emergency medical care.}{ref72}
\indianentry{IND-070}{A dog bite can be treated with chilli powder, turmeric, oil, or another traditional application.}{False}{Household substances do not prevent rabies or reliably disinfect a bite. Wash the wound immediately with soap and running water and obtain urgent assessment for rabies prophylaxis, tetanus, antibiotics, and wound care. Rabies is preventable after exposure but almost always fatal after symptoms begin.}{ref73}
\indianentry{IND-071}{A person bitten by a dog should avoid particular foods while receiving rabies vaccination.}{False}{Standard rabies vaccination does not require traditional food restrictions. A normal balanced diet is usually appropriate unless a clinician identifies another issue. The critical steps are wound washing, timely vaccination, and rabies immunoglobulin when indicated.}{ref73}
\indianentry{IND-072}{The old ``14 injections in the stomach'' are still required after dog bites.}{False}{Modern rabies post-exposure regimens do not use the historical abdominal injection schedule. The schedule and route depend on the vaccine, exposure, prior vaccination, immune status, and national guidance. Some exposures also require rabies immunoglobulin.}{ref73}
\indianentry{IND-073}{A child with measles should not be bathed because the rash will go back inside the body.}{False}{Bathing does not force a measles rash inside the body. Gentle washing can support comfort and hygiene if the child is stable and kept warm. Measles can cause pneumonia, dehydration, encephalitis, and other complications, so medical and public-health advice matters.}{ref75}
\indianentry{IND-074}{Measles is a visitation of Sheetala Mata and should not be medically treated.}{False}{Measles is caused by measles virus and can produce serious complications. Cultural interpretation may be meaningful, but it does not replace diagnosis, supportive care, vitamin A when recommended, or treatment of complications. Breathing difficulty, dehydration, altered consciousness, or persistent fever requires urgent care.}{ref75}
\indianentry{IND-075}{Chickenpox is simply a manifestation of a goddess, so medicines should always be avoided.}{False}{Chickenpox is caused by varicella-zoster virus. Many cases are self-limited, but complications occur particularly in adults, pregnancy, infants, and people with weakened immunity. Medical assessment may identify a need for antiviral treatment, symptom control, or urgent care.}{ref76}
\indianentry{IND-076}{Smallpox was caused by divine displeasure.}{False}{Smallpox was caused by variola virus, not divine displeasure. It was eradicated through surveillance, isolation, and vaccination. The history demonstrates the value of evidence-based public health for infectious disease.}{ref77}
\indianentry{IND-077}{Food restrictions alone can cure jaundice.}{False}{Jaundice is a sign of excess bilirubin, not a single disease. Causes include hepatitis, medication injury, obstruction, blood disorders, and liver disease. Treatment depends on the cause and may require tests, medicines, procedures, or hospital care.}{ref78}
\indianentry{IND-078}{A person with jaundice should eat only bland or yellow-free food and avoid normal nutrition.}{False}{Food color does not determine bilirubin or cure jaundice. Unless a clinician gives a specific restriction, people generally need adequate energy, protein, and fluids appropriate to their condition. Yellow eyes or skin, dark urine, pale stool, fever, pain, confusion, or bleeding requires assessment.}{ref78}
\indianentry{IND-079}{Sugarcane juice cures jaundice.}{False}{Sugarcane juice does not treat the diseases that can cause jaundice. Unpasteurized juice may carry contamination, and large amounts can be harmful for people with diabetes or vomiting. It is not a substitute for diagnosis or treatment.}{ref78}
\indianentry{IND-080}{A traditional healer can remove jaundice by transferring it into water, leaves, or another object.}{False}{Changing an object cannot measure or remove bilirubin from the body. Temporary symptom changes do not establish that liver or blood disease has resolved. Jaundice should be assessed by a qualified clinician with appropriate examination and testing.}{ref78}
\end{enumerate}
